\section{Data and models}
\label{sec:data}

\paragraph{Tier A: public benchmark.} Thirty of the 100 fev-bench tasks \citep{shchur2025fevbench}, sampled with a fixed seed and
stratified by domain and frequency, after excluding the three tasks whose string covariates Google's runner encodes with Python's
salted \texttt{hash()} (not reproducible across processes). Metrics are computed with fev's own functions, and every Tier A model uses
the conventions of Google's runner (last-value imputation, context cut to 15{,}360 points). Two pipeline controls are required on
every task: our Seasonal Naive must match fev's published results ($|\Delta\mathrm{MASE}| < 5\cdot10^{-5}$) and our Chronos-2 fev's
published Chronos-2 ($|\Delta\mathrm{SQL}| \le 0.5\%$).

\paragraph{Tier B: post-release real data.} TimesFM~3 was released on 2026-08-31. The test windows run from 2026-09-01 to
2026-10-20, with origins every 24 hours for 24-hour horizons and weekly otherwise (168 hours, 7 days or 5 business days). Six domains, with
series lists fixed by rules on 2025 data: hourly electricity load in 12 European countries and day-ahead prices in 12 bidding zones
(Energy-Charts \citep{energycharts2026}), hourly temperature and wind speed in 20 cities (ERA5 via Open-Meteo
\citep{hersbach2020era5,zippenfenig2023openmeteo}), daily flow at 20 river gauges (USGS \citep{usgs2026nwis}), ten business-daily exchange rates
against the euro (ECB \citep{ecb2026dataportal}), and eight Treasury yields with WTI and Brent prices (FRED \citep{fred2026}). Horizons are 24 and
168 hours, 7 and 28 days, 5 and 20 business days. Load and river flow have a covariate whose \emph{observed} future values are
supplied (temperature, precipitation): it is an oracle, used only in secondary analyses as an upper bound. Price, FX and FRED also
form one group task each, in which all series of the domain are forecast jointly. Same-source data of earlier periods are part of
the public pre-training corpora for load, price and weather; the model card describes the TimesFM~3 corpus only at the level of sources (GiftEvalPretrain without the
datasets that overlap fev-bench, Wikipedia pageviews to November 2023, Google Trends to 2022, synthetic data
\citep{timesfm2026hf}), so Tier B guarantees
that the \emph{test windows} are new, not that the series are.

\paragraph{Tier W: memorisation probe.} English Wikipedia daily user pageviews \citep{wikimedia2026pageviews}, horizons 7 and 28 days,
400-day contexts. \emph{W-seen}: 200 popular pages (seed 17) with test windows from 2023-09-01 to 2023-10-20, inside the pageview
corpus that the TimesFM~3 model card lists up to November 2023 \citep{timesfm2026hf}. \emph{W-post}: the same pages, same calendar span in
2026. \emph{W-new}: pages that first entered the top-1000 lists in 2024--2025, with no views at all before December 2023, matched to
W-seen/W-post on quintiles of median views. Pages that became redirects are excluded (190 and 128 pages remain).

\paragraph{Tier C: controlled synthetic data.} 38 cells with 200 replications each (7{,}600 tasks), seeded per cell and
replication: (i) a replication of the three hard processes of a recent simulation study of zero-shot forecasting
\citep{ebrahim2026foundation}, level shift (D6), logistic growth (D7) and intermittent demand (D8), with contexts of 48 and 200 points and
horizons 12 and 48; (ii) univariate stress cells: two seasonalities with one or four weeks of context, a season longer than the
context, a context ladder from 64 to 15{,}360 points, horizons up to 1{,}024, a regime change near the origin, outliers in the last
5\% of the context, a random walk and white noise; (iii) multivariate cells: VAR(1) systems of 2 to 64 variates, lead-lag pairs
(lags 1, 5, 24), a target driven by a known future covariate (signal-to-noise 0.5 and 2), a spurious past covariate, and independent
series stacked as variates.

\paragraph{Tier P: prospective check.} On the morning of each issue day (2026-10-30 and 2026-11-02) the data published up to the
previous day are downloaded for the load, price, river, FX, FRED and Wikipedia series, every series is forecast from its last
observation, and the contexts, tasks, code bundle, run records and forecast files are hashed together; the combined SHA-256 is
recorded before the scoring window (from two to nine days after the issue day) starts.

\paragraph{Models.} Table~\ref{tab:models} lists the configurations. Every Hugging Face model is loaded by revision hash.
TimesFM~3 runs in four configurations: package defaults in multivariate mode with covariates (\texttt{tfm3}), univariate without
covariates (\texttt{tfm3-uni}, the configuration of the information-matched contrasts), the flags of Google's benchmark runners
(\texttt{tfm3-eval}), and univariate with the classical models' context of 2{,}016 points (\texttt{tfm3-c2016}). The comparison
foundation models are TimesFM~2.5 \citep{das2024timesfm}, Chronos-2 \citep{ansari2025chronos2} and TiRex-2 \citep{podest2026tirex2}, each also
in a univariate configuration where relevant. The classical models are StatsForecast \citep{garza2022statsforecast} implementations of
AutoETS, AutoARIMA and AutoTheta, MSTL with seasons 24 and 168 \citep{bandara2025mstl} for hourly data, and Naive, Seasonal Naive and
the historic mean. Classical models use at most the last 2{,}016 hourly, 1{,}095 daily or 780 business-daily observations,
the foundation models up to 15{,}360 (Chronos-2 truncates at 8{,}192). On daily and business-daily data all models therefore see the
same history; on hourly data the foundation models see about 15{,}000 points, and \texttt{tfm3-c2016} measures the effect of that
difference for TimesFM~3. The statistical ensemble averages the quantiles of three members: MSTL with ETS, Theta and ARIMA remainders on
hourly data, and AutoETS, AutoARIMA and AutoTheta otherwise. All foundation models ran on a free Google Colab Tesla T4 GPU in fp32;
classical models ran on eight cores of a local CPU.

\begin{table}[!htbp]
\centering
\caption{Model configurations.}
\label{tab:models}
\scriptsize
\begin{tabular}{llp{7.6cm}}
\toprule
Key & Model (revision) & Configuration \\
\midrule
\texttt{tfm3} & TimesFM~3, 330M (\texttt{43046b85}) & package defaults; all targets and covariates of a task in one pass \\
\texttt{tfm3-uni} & TimesFM~3 & package defaults; each target alone, no covariates \\
\texttt{tfm3-eval} & TimesFM~3 & as \texttt{tfm3} with symmetric averaging and \emph{make positive} \\
\texttt{tfm3-c2016} & TimesFM~3 & as \texttt{tfm3-uni}, last 2{,}016 points only \\
\texttt{tfm25} & TimesFM~2.5, 200M (\texttt{1d952420}) & univariate, recommended flags \\
\texttt{chronos2}, \texttt{-uni} & Chronos-2, 120M (\texttt{29ec3766}) & \texttt{predict\_df}; with / without covariates and joint targets \\
\texttt{tirex2}, \texttt{-uni} & TiRex-2 (\texttt{30788a31}) & as Chronos-2; horizons up to 320 \\
\texttt{stat-ens} & equal-weight ensemble & MSTL-\{ETS, Theta, ARIMA\} (hourly); AutoETS, AutoARIMA, AutoTheta \\
classical & StatsForecast 2.1.1 & Naive, Seasonal Naive, mean, AutoETS, AutoARIMA, AutoTheta, MSTL variants \\
\bottomrule
\end{tabular}
\end{table}
